\section{Results}
\label{sec:results}

All numbers in this section come from \resNFits{} outer-fold fits
(\resNSeeds{} seeds $\times$ 5 folds $\times$ 6 datasets $\times$ 11 models). They are
summarised as the mean and standard deviation over the \resFoldsPerDataset{} test folds
per dataset. Tables~\ref{tab:auroc}--\ref{tab:ece} report discrimination and
calibration, Table~\ref{tab:hyp} the preregistered hypotheses, and
Table~\ref{tab:complexity} model size and cost. PR-AUC and calibration slopes are in
the appendix (Tables~\ref{tab:prauc} and~\ref{tab:slope}).

\begin{table*}[t]
\centering
\scriptsize
\setlength{\tabcolsep}{3pt}
\caption{AUROC (higher is better). Mean\,$\pm$\,SD over 25 outer test folds (5 seeds $\times$ 5 folds); bold = best of the 11 preregistered models per dataset. $^\dagger$Ablation from the same NSGA-II runs (not ranked).}
\label{tab:auroc}
\begin{adjustbox}{max width=\linewidth}
\begin{tabular}{lcccccc}
\toprule
 & \multicolumn{2}{c}{Development} & \multicolumn{4}{c}{Held-out} \\
\cmidrule(lr){2-3}\cmidrule(lr){4-7}
Model & Pima & Cleveland & WDBC & ILPD & Thyroid & CKD \\
\midrule
LogReg & 0.834\,$\pm$\,0.037 & \textbf{0.912\,$\pm$\,0.031} & \textbf{0.996\,$\pm$\,0.005} & \textbf{0.751\,$\pm$\,0.033} & 0.953\,$\pm$\,0.017 & 1.000\,$\pm$\,0.000 \\
Decision tree & 0.758\,$\pm$\,0.045 & 0.798\,$\pm$\,0.046 & 0.953\,$\pm$\,0.026 & 0.695\,$\pm$\,0.042 & 0.974\,$\pm$\,0.020 & 0.975\,$\pm$\,0.017 \\
XGBoost & 0.837\,$\pm$\,0.038 & 0.904\,$\pm$\,0.028 & 0.993\,$\pm$\,0.006 & 0.726\,$\pm$\,0.032 & 0.997\,$\pm$\,0.002 & 1.000\,$\pm$\,0.000 \\
LightGBM & 0.836\,$\pm$\,0.034 & 0.903\,$\pm$\,0.029 & 0.994\,$\pm$\,0.005 & 0.721\,$\pm$\,0.033 & 0.998\,$\pm$\,0.001 & \textbf{1.000\,$\pm$\,0.000} \\
CatBoost & 0.839\,$\pm$\,0.036 & 0.907\,$\pm$\,0.028 & 0.994\,$\pm$\,0.006 & 0.743\,$\pm$\,0.036 & \textbf{0.998\,$\pm$\,0.002} & 1.000\,$\pm$\,0.000 \\
MLP & 0.827\,$\pm$\,0.037 & 0.889\,$\pm$\,0.048 & 0.991\,$\pm$\,0.010 & 0.686\,$\pm$\,0.081 & 0.973\,$\pm$\,0.018 & 0.996\,$\pm$\,0.007 \\
XGBoost + Platt & 0.839\,$\pm$\,0.038 & 0.906\,$\pm$\,0.030 & 0.992\,$\pm$\,0.006 & 0.728\,$\pm$\,0.036 & 0.996\,$\pm$\,0.003 & 1.000\,$\pm$\,0.001 \\
XGBoost + isotonic & 0.837\,$\pm$\,0.040 & 0.905\,$\pm$\,0.030 & 0.991\,$\pm$\,0.008 & 0.725\,$\pm$\,0.034 & 0.994\,$\pm$\,0.007 & 0.999\,$\pm$\,0.003 \\
SO-GA-Fuzzy (classic) & 0.793\,$\pm$\,0.036 & 0.890\,$\pm$\,0.037 & 0.987\,$\pm$\,0.010 & 0.636\,$\pm$\,0.067 & 0.845\,$\pm$\,0.141 & 0.997\,$\pm$\,0.008 \\
SO-GA-Fuzzy (hybrid) & 0.820\,$\pm$\,0.033 & 0.878\,$\pm$\,0.033 & 0.992\,$\pm$\,0.007 & 0.690\,$\pm$\,0.058 & 0.987\,$\pm$\,0.011 & 0.998\,$\pm$\,0.005 \\
MO-Fuzzy (NSGA-II) & \textbf{0.841\,$\pm$\,0.038} & 0.892\,$\pm$\,0.040 & 0.988\,$\pm$\,0.009 & 0.710\,$\pm$\,0.045 & 0.976\,$\pm$\,0.008 & 0.993\,$\pm$\,0.007 \\
\midrule
MO-Fuzzy (best-Brier)$^\dagger$ & 0.837\,$\pm$\,0.039 & 0.896\,$\pm$\,0.033 & 0.988\,$\pm$\,0.010 & 0.712\,$\pm$\,0.046 & 0.960\,$\pm$\,0.023 & 0.999\,$\pm$\,0.002 \\
\bottomrule
\end{tabular}
\end{adjustbox}
\end{table*}

\begin{table*}[t]
\centering
\scriptsize
\setlength{\tabcolsep}{3pt}
\caption{Brier score (lower is better). Mean\,$\pm$\,SD over 25 outer test folds (5 seeds $\times$ 5 folds); bold = best of the 11 preregistered models per dataset. $^\dagger$Ablation from the same NSGA-II runs (not ranked).}
\label{tab:brier}
\begin{adjustbox}{max width=\linewidth}
\begin{tabular}{lcccccc}
\toprule
 & \multicolumn{2}{c}{Development} & \multicolumn{4}{c}{Held-out} \\
\cmidrule(lr){2-3}\cmidrule(lr){4-7}
Model & Pima & Cleveland & WDBC & ILPD & Thyroid & CKD \\
\midrule
LogReg & 0.158\,$\pm$\,0.017 & \textbf{0.115\,$\pm$\,0.021} & \textbf{0.020\,$\pm$\,0.007} & \textbf{0.172\,$\pm$\,0.008} & 0.027\,$\pm$\,0.003 & 0.005\,$\pm$\,0.005 \\
Decision tree & 0.181\,$\pm$\,0.017 & 0.179\,$\pm$\,0.029 & 0.061\,$\pm$\,0.019 & 0.185\,$\pm$\,0.010 & 0.015\,$\pm$\,0.003 & 0.036\,$\pm$\,0.020 \\
XGBoost & 0.156\,$\pm$\,0.016 & 0.124\,$\pm$\,0.022 & 0.025\,$\pm$\,0.008 & 0.180\,$\pm$\,0.008 & 0.009\,$\pm$\,0.003 & 0.006\,$\pm$\,0.005 \\
LightGBM & 0.157\,$\pm$\,0.016 & 0.123\,$\pm$\,0.022 & 0.024\,$\pm$\,0.008 & 0.182\,$\pm$\,0.008 & 0.008\,$\pm$\,0.002 & \textbf{0.002\,$\pm$\,0.002} \\
CatBoost & 0.156\,$\pm$\,0.016 & 0.122\,$\pm$\,0.021 & 0.023\,$\pm$\,0.009 & 0.174\,$\pm$\,0.008 & \textbf{0.007\,$\pm$\,0.002} & 0.005\,$\pm$\,0.005 \\
MLP & 0.163\,$\pm$\,0.016 & 0.137\,$\pm$\,0.034 & 0.032\,$\pm$\,0.020 & 0.185\,$\pm$\,0.014 & 0.021\,$\pm$\,0.004 & 0.048\,$\pm$\,0.031 \\
XGBoost + Platt & 0.156\,$\pm$\,0.017 & 0.122\,$\pm$\,0.020 & 0.026\,$\pm$\,0.009 & 0.180\,$\pm$\,0.009 & 0.012\,$\pm$\,0.003 & 0.008\,$\pm$\,0.006 \\
XGBoost + isotonic & 0.156\,$\pm$\,0.018 & 0.122\,$\pm$\,0.022 & 0.026\,$\pm$\,0.008 & 0.180\,$\pm$\,0.008 & 0.011\,$\pm$\,0.003 & 0.008\,$\pm$\,0.007 \\
SO-GA-Fuzzy (classic) & 0.178\,$\pm$\,0.011 & 0.133\,$\pm$\,0.025 & 0.048\,$\pm$\,0.019 & 0.202\,$\pm$\,0.017 & 0.080\,$\pm$\,0.039 & 0.067\,$\pm$\,0.021 \\
SO-GA-Fuzzy (hybrid) & 0.167\,$\pm$\,0.018 & 0.151\,$\pm$\,0.031 & 0.033\,$\pm$\,0.014 & 0.197\,$\pm$\,0.019 & 0.017\,$\pm$\,0.004 & 0.011\,$\pm$\,0.008 \\
MO-Fuzzy (NSGA-II) & \textbf{0.155\,$\pm$\,0.017} & 0.132\,$\pm$\,0.028 & 0.036\,$\pm$\,0.013 & 0.183\,$\pm$\,0.012 & 0.023\,$\pm$\,0.004 & 0.021\,$\pm$\,0.011 \\
\midrule
MO-Fuzzy (best-Brier)$^\dagger$ & 0.156\,$\pm$\,0.018 & 0.129\,$\pm$\,0.025 & 0.035\,$\pm$\,0.014 & 0.184\,$\pm$\,0.013 & 0.017\,$\pm$\,0.004 & 0.011\,$\pm$\,0.010 \\
\bottomrule
\end{tabular}
\end{adjustbox}
\end{table*}

\begin{table*}[t]
\centering
\scriptsize
\setlength{\tabcolsep}{3pt}
\caption{Expected calibration error, 10 equal-width bins (lower is better). Mean\,$\pm$\,SD over 25 outer test folds (5 seeds $\times$ 5 folds); bold = best of the 11 preregistered models per dataset. $^\dagger$Ablation from the same NSGA-II runs (not ranked).}
\label{tab:ece}
\begin{adjustbox}{max width=\linewidth}
\begin{tabular}{lcccccc}
\toprule
 & \multicolumn{2}{c}{Development} & \multicolumn{4}{c}{Held-out} \\
\cmidrule(lr){2-3}\cmidrule(lr){4-7}
Model & Pima & Cleveland & WDBC & ILPD & Thyroid & CKD \\
\midrule
LogReg & 0.076\,$\pm$\,0.014 & 0.104\,$\pm$\,0.023 & 0.034\,$\pm$\,0.009 & 0.072\,$\pm$\,0.022 & 0.021\,$\pm$\,0.004 & 0.013\,$\pm$\,0.009 \\
Decision tree & \textbf{0.059\,$\pm$\,0.025} & \textbf{0.088\,$\pm$\,0.032} & 0.039\,$\pm$\,0.018 & \textbf{0.057\,$\pm$\,0.029} & \textbf{0.007\,$\pm$\,0.003} & 0.026\,$\pm$\,0.013 \\
XGBoost & 0.082\,$\pm$\,0.019 & 0.113\,$\pm$\,0.028 & 0.032\,$\pm$\,0.009 & 0.081\,$\pm$\,0.020 & 0.009\,$\pm$\,0.003 & 0.023\,$\pm$\,0.009 \\
LightGBM & 0.084\,$\pm$\,0.018 & 0.107\,$\pm$\,0.027 & \textbf{0.029\,$\pm$\,0.008} & 0.079\,$\pm$\,0.022 & 0.008\,$\pm$\,0.003 & \textbf{0.005\,$\pm$\,0.005} \\
CatBoost & 0.075\,$\pm$\,0.023 & 0.111\,$\pm$\,0.025 & 0.033\,$\pm$\,0.008 & 0.073\,$\pm$\,0.020 & 0.007\,$\pm$\,0.002 & 0.013\,$\pm$\,0.008 \\
MLP & 0.083\,$\pm$\,0.021 & 0.120\,$\pm$\,0.031 & 0.050\,$\pm$\,0.032 & 0.072\,$\pm$\,0.026 & 0.017\,$\pm$\,0.005 & 0.126\,$\pm$\,0.068 \\
XGBoost + Platt & 0.083\,$\pm$\,0.019 & 0.113\,$\pm$\,0.025 & 0.041\,$\pm$\,0.011 & 0.083\,$\pm$\,0.026 & 0.011\,$\pm$\,0.003 & 0.031\,$\pm$\,0.009 \\
XGBoost + isotonic & 0.081\,$\pm$\,0.018 & 0.103\,$\pm$\,0.033 & 0.033\,$\pm$\,0.012 & 0.076\,$\pm$\,0.022 & 0.009\,$\pm$\,0.002 & 0.015\,$\pm$\,0.009 \\
SO-GA-Fuzzy (classic) & 0.104\,$\pm$\,0.027 & 0.125\,$\pm$\,0.029 & 0.106\,$\pm$\,0.041 & 0.105\,$\pm$\,0.042 & 0.185\,$\pm$\,0.095 & 0.186\,$\pm$\,0.040 \\
SO-GA-Fuzzy (hybrid) & 0.091\,$\pm$\,0.032 & 0.142\,$\pm$\,0.030 & 0.038\,$\pm$\,0.015 & 0.118\,$\pm$\,0.033 & 0.013\,$\pm$\,0.003 & 0.020\,$\pm$\,0.013 \\
MO-Fuzzy (NSGA-II) & 0.080\,$\pm$\,0.015 & 0.109\,$\pm$\,0.025 & 0.040\,$\pm$\,0.010 & 0.083\,$\pm$\,0.029 & 0.013\,$\pm$\,0.004 & 0.023\,$\pm$\,0.009 \\
\midrule
MO-Fuzzy (best-Brier)$^\dagger$ & 0.081\,$\pm$\,0.023 & 0.119\,$\pm$\,0.020 & 0.041\,$\pm$\,0.015 & 0.093\,$\pm$\,0.029 & 0.011\,$\pm$\,0.004 & 0.017\,$\pm$\,0.011 \\
\bottomrule
\end{tabular}
\end{adjustbox}
\end{table*}


\subsection{Discrimination and calibration}
Averaged over the six datasets, MO-Fuzzy reached an AUROC of
\resMeanAurocAllMOFuzzyNSGAII{} and a Brier score of \resMeanBrierAllMOFuzzyNSGAII{}.
CatBoost, the strongest model, reached \resMeanAurocAllCatBoost{} and
\resMeanBrierAllCatBoost{}, and L2 logistic regression reached \resMeanAurocAllLogReg{}
and \resMeanBrierAllLogReg{}. The classic genetic-fuzzy baseline reached
\resMeanAurocAllSOGAFuzzyClassic{} and \resMeanBrierAllSOGAFuzzyClassic{}. On the four
held-out datasets the mean Brier scores were \resMeanBrierHeldoutMOFuzzyNSGAII{}
(MO-Fuzzy), \resMeanBrierHeldoutSOGAFuzzyHybrid{} (SO-GA-Fuzzy hybrid),
\resMeanBrierHeldoutSOGAFuzzyClassic{} (SO-GA-Fuzzy classic), \resMeanBrierHeldoutCatBoost{}
(CatBoost) and \resMeanBrierHeldoutLogReg{} (logistic regression).

The ranking differs by dataset. On Pima, MO-Fuzzy had the highest AUROC
(\resMOAurocPima{}) and the lowest Brier score (\resBestBrierPima{}) of all eleven
models. On Cleveland, WDBC and ILPD the best model in both metrics was logistic
regression, and MO-Fuzzy trailed it by up to a few hundredths. On the thyroid and CKD
data, which are close to separable, the best tree ensembles reached near-perfect
discrimination and Brier scores of \resBestBrierThyroid{} and \resBestBrierCKD{}. There
MO-Fuzzy, with on average \resMORulesThyroid{} and \resMORulesCKD{} rules, discriminated
well (AUROC \resMOAurocThyroid{} and \resMOAurocCKD{}) but was clearly less sharp than
the ensembles (Brier \resMOBrierThyroid{} and \resMOBrierCKD{}).

Among the genetic-fuzzy models, MO-Fuzzy had the lowest mean ECE
(\resMeanEceAllMOFuzzyNSGAII{}, against \resMeanEceAllSOGAFuzzyHybrid{} for the hybrid
and \resMeanEceAllSOGAFuzzyClassic{} for the classic baseline). This is close to the
tuned ensembles (\resMeanEceAllCatBoost{} for CatBoost and \resMeanEceAllXGBoost{} for
XGBoost). MO-Fuzzy had a lower Brier score than the single-objective hybrid on
\resNBrierMOBetterHybrid{} of 6 datasets and an AUROC at least as high on
\resNAurocMOAtLeastHybrid{} of 6, with a fraction of its rules
(Section~\ref{sec:complexity}). Its median absolute deviation of the calibration slope
from 1 was between \resMOSlopeMedianDevThyroid{} (thyroid) and
\resMOSlopeMedianDevCKD{} (CKD); Table~\ref{tab:slope} gives the slopes. On CKD the
slope is ill-determined for every model because the predictions separate the classes
almost perfectly. The pooled reliability diagrams (Fig.~\ref{fig:reliability}) show
that MO-Fuzzy probabilities follow the diagonal about as closely as those of XGBoost and
CatBoost. The classic genetic-fuzzy model is systematically miscalibrated on the
held-out data: its evolved consequents give step-like probabilities that over- or
under-shoot whole risk bands.

\begin{figure}[t]
\centering
\includegraphics[width=\linewidth]{reliability}
\caption{Reliability diagrams pooled over all \resFoldsPerDataset{} outer test folds
(10 equal-width bins; bins with fewer than 10 pooled predictions are omitted). The dashed diagonal is
perfect calibration.}
\label{fig:reliability}
\end{figure}

\subsection{Preregistered hypotheses}
\begin{table}[t]
\centering
\scriptsize
\setlength{\tabcolsep}{3pt}
\caption{Preregistered hypotheses (PREREGISTRATION.md \S5) and outcomes.}
\label{tab:hyp}
\begin{adjustbox}{max width=\linewidth}
\begin{tabular}{lp{5.4cm}lp{1.6cm}}
\toprule
ID & Criterion & Observed & Outcome \\
\midrule
H1 & Brier: MO-Fuzzy $<$ SO-GA-Fuzzy (classic) on $\geq$5/6 datasets and one-sided Wilcoxon $p<0.05$ & 6/6, $p=0.016$ & \textbf{Supported} \\
H2 & AUROC: MO-Fuzzy $\geq$ best GBDT $-0.02$ on $\geq$4/6 datasets & 4/6 & \textbf{Supported} \\
H3 & $|$slope$-1|$: MO-Fuzzy $<$ SO-GA-Fuzzy (classic) on $\geq$4/6 and one-sided Wilcoxon $p<0.05$ & 6/6, $p=0.016$ & \textbf{Supported} \\
H4 & Literals: MO $\leq 0.5\times$ SO-hybrid on 6/6 and AUROC loss $\leq 0.01$ on $\geq$4/6 & 6/6 and 5/6 & \textbf{Supported} \\
\bottomrule
\end{tabular}
\end{adjustbox}
\end{table}

Table~\ref{tab:hyp} gives the outcome of each preregistered criterion.
\textbf{H1} (Brier: MO-Fuzzy better than the classic genetic-fuzzy model) was
\resHOne{}: MO-Fuzzy had the lower Brier score on \resHOneWins{}/6 datasets
(one-sided Wilcoxon $p=\resHOneP{}$, the smallest attainable with $N=6$).
\textbf{H2} (AUROC within 0.02 of the best gradient-boosted tree on at least four
datasets) was \resHTwo{} with exactly \resHTwoN{}/6. It failed on ILPD (difference
\resGbdtGapILPD{}) and thyroid (\resGbdtGapThyroid{}), so the margin is narrow.
\textbf{H3} (calibration slope closer to 1 than the classic model) was \resHThree{}
(\resHThreeWins{}/6, $p=\resHThreeP{}$). \textbf{H4} (at most half the literals of the
single-objective hybrid with an AUROC loss of at most 0.01 on four or more datasets) was
\resHFour{}. The literal ratio ranged from \resHFourRatioMin{} to \resHFourRatioMax{}
on the six datasets, and the AUROC loss was within 0.01 on \resHFourLossN{}/6.

\subsection{Across-dataset comparison}
\begin{table}[t]
\centering
\scriptsize
\setlength{\tabcolsep}{4pt}
\caption{Average ranks over the 6 datasets (1 = best) and two-sided exact Wilcoxon signed-rank tests of MO-Fuzzy against each model on dataset means, Holm-corrected. With $N=6$ the smallest attainable two-sided $p$ is 0.031, so no comparison can survive Holm correction (declared in the preregistration).}
\label{tab:ranks}
\begin{adjustbox}{max width=\linewidth}
\begin{tabular}{lrrrrrr}
\toprule
 & \multicolumn{3}{c}{AUROC} & \multicolumn{3}{c}{Brier} \\
\cmidrule(lr){2-4}\cmidrule(lr){5-7}
Model & Rank & $p$ & $p_{\mathrm{Holm}}$ & Rank & $p$ & $p_{\mathrm{Holm}}$ \\
\midrule
LogReg & 3.75 & 0.438 & 1.000 & 3.83 & 0.156 & 0.562 \\
Decision tree & 10.00 & 0.031 & 0.312 & 9.33 & 0.094 & 0.562 \\
XGBoost & 4.17 & 0.062 & 0.562 & 4.17 & 0.062 & 0.562 \\
LightGBM & 4.00 & 0.062 & 0.562 & 3.83 & 0.094 & 0.562 \\
CatBoost & 2.08 & 0.062 & 0.562 & 2.17 & 0.062 & 0.562 \\
MLP & 8.67 & 0.438 & 1.000 & 8.50 & 0.219 & 0.562 \\
XGBoost + Platt & 3.67 & 0.062 & 0.562 & 4.67 & 0.062 & 0.562 \\
XGBoost + isotonic & 5.33 & 0.094 & 0.562 & 4.00 & 0.062 & 0.562 \\
SO-GA-Fuzzy (classic) & 9.67 & 0.156 & 0.625 & 10.17 & 0.031 & 0.312 \\
SO-GA-Fuzzy (hybrid) & 7.83 & 0.438 & 1.000 & 8.50 & 0.438 & 0.562 \\
MO-Fuzzy (NSGA-II) & 6.83 & -- & -- & 6.83 & -- & -- \\
\midrule
Friedman $\chi^2_{10}$ & \multicolumn{3}{c}{40.54 ($p=$<$0.001$)} & \multicolumn{3}{c}{40.85 ($p=$<$0.001$)} \\
Nemenyi CD ($\alpha=0.05$) & \multicolumn{3}{c}{6.16} & \multicolumn{3}{c}{6.16} \\
\bottomrule
\end{tabular}
\end{adjustbox}
\end{table}

The Friedman test rejected equal performance for both AUROC
($\chi^2_{10}=\resFriedmanAurocChi{}$, $p=\resFriedmanAurocP{}$) and Brier score
($\chi^2_{10}=\resFriedmanBrierChi{}$, $p=\resFriedmanBrierP{}$). CatBoost had the best
average rank (AUROC \resBestRankAuroc{}, Brier \resBestRankBrier{}). MO-Fuzzy ranked
\resRankMOAuroc{} and \resRankMOBrier{}, the best of the three genetic-fuzzy models in
both metrics (Fig.~\ref{fig:cd}). Its rank gap to CatBoost (\resRankGapAuroc{} for AUROC
and \resRankGapBrier{} for Brier) is \resRankGapVsCdAuroc{} and \resRankGapVsCdBrier{},
respectively, than the Nemenyi critical difference of \resCdAuroc{}. As the
preregistration anticipated, no Holm-corrected Wilcoxon comparison of MO-Fuzzy was
significant (smallest $p_{\mathrm{Holm}}$: \resMinWilcoxonPHolmAuroc{} for AUROC,
\resMinWilcoxonPHolmBrier{} for Brier). With six datasets these tests cannot separate
models whose ranks differ by less than the critical difference. The rankings are
therefore descriptive.

\begin{figure}[t]
\centering
\includegraphics[width=\linewidth]{cd_auroc}\\[2pt]
\includegraphics[width=\linewidth]{cd_brier}
\caption{Critical-difference diagrams of average ranks over the six datasets
(Nemenyi, $\alpha=0.05$; top: AUROC, bottom: Brier). Models joined by a bar are not
significantly different.}
\label{fig:cd}
\end{figure}

\begin{table*}[t]
\centering
\scriptsize
\setlength{\tabcolsep}{3pt}
\caption{Per-dataset difference MO-Fuzzy $-$ model (AUROC / Brier) with Nadeau--Bengio corrected resampled $t$-tests over 25 paired folds, Holm-corrected within dataset and metric. $^{*}$: $p_{\mathrm{Holm}}<0.05$. Positive AUROC and negative Brier differences favour MO-Fuzzy.}
\label{tab:nb}
\begin{adjustbox}{max width=\linewidth}
\begin{tabular}{lcccccc}
\toprule
Model & Pima & Cleveland & WDBC & ILPD & Thyroid & CKD \\
\midrule
LogReg & +0.007 / -0.002 & -0.020 / +0.016 & -0.008 / +0.016$^{*}$ & -0.040 / +0.011 & +0.023 / -0.004 & -0.007 / +0.016$^{*}$ \\
Decision tree & +0.082$^{*}$ / -0.026 & +0.094$^{*}$ / -0.048 & +0.035 / -0.025$^{*}$ & +0.016 / -0.002 & +0.002 / +0.007$^{*}$ & +0.017 / -0.015 \\
XGBoost & +0.004 / -0.001 & -0.012 / +0.008 & -0.005 / +0.011 & -0.016 / +0.003 & -0.021$^{*}$ / +0.013$^{*}$ & -0.007 / +0.016$^{*}$ \\
LightGBM & +0.005 / -0.002 & -0.011 / +0.009 & -0.006 / +0.012 & -0.011 / +0.001 & -0.022$^{*}$ / +0.015$^{*}$ & -0.007 / +0.020$^{*}$ \\
CatBoost & +0.002 / -0.000 & -0.015 / +0.010 & -0.006 / +0.013$^{*}$ & -0.032 / +0.009 & -0.022$^{*}$ / +0.016$^{*}$ & -0.007 / +0.016$^{*}$ \\
MLP & +0.014 / -0.007 & +0.002 / -0.005 & -0.003 / +0.004 & +0.025 / -0.002 & +0.003 / +0.001 & -0.004 / -0.026 \\
XGBoost + Platt & +0.001 / -0.001 & -0.014 / +0.010 & -0.005 / +0.010 & -0.017 / +0.003 & -0.020$^{*}$ / +0.011$^{*}$ & -0.007 / +0.013$^{*}$ \\
XGBoost + isotonic & +0.003 / -0.000 & -0.014 / +0.010 & -0.003 / +0.010 & -0.015 / +0.004 & -0.018$^{*}$ / +0.012$^{*}$ & -0.006 / +0.014$^{*}$ \\
SO-GA-Fuzzy (classic) & +0.047 / -0.023 & +0.002 / -0.002 & +0.001 / -0.012 & +0.074 / -0.019 & +0.130 / -0.057$^{*}$ & -0.004 / -0.046$^{*}$ \\
SO-GA-Fuzzy (hybrid) & +0.021 / -0.011 & +0.014 / -0.019 & -0.004 / +0.003 & +0.021 / -0.014 & -0.011 / +0.006$^{*}$ & -0.006 / +0.010 \\
\bottomrule
\end{tabular}
\end{adjustbox}
\end{table*}

The per-dataset corrected resampled $t$-tests (Table~\ref{tab:nb}) locate the
differences. On Pima, Cleveland and ILPD, no difference between MO-Fuzzy and any
gradient-boosted, calibrated or logistic model was significant, and MO-Fuzzy was
significantly better than the decision tree in AUROC on Pima and Cleveland. The
significant deficits of MO-Fuzzy are concentrated on the three near-separable held-out
datasets. On thyroid it was worse than all gradient-boosted variants in both metrics,
and in Brier score also worse than the decision tree and the single-objective hybrid;
on CKD it had a worse Brier score than logistic regression and all gradient-boosted
variants; and on WDBC it had a worse Brier score than logistic regression and CatBoost.
Against the classic genetic-fuzzy model, MO-Fuzzy had a significantly better Brier score
on thyroid and CKD.

\subsection{Pareto fronts and the choice of operating point}
\label{sec:pareto}
\begin{figure}[t]
\centering
\includegraphics[width=\linewidth]{pareto_fronts}
\caption{Test-fold performance of all external-archive members of MO-Fuzzy against
their total number of literals (grey points: members from all folds; line: median per
literal count). Stars mark the preregistered $\varepsilon$-compact operating point and
diamonds the best-Brier ablation (means over folds). The dashed line is the best
gradient-boosted tree.}
\label{fig:pareto}
\end{figure}
\begin{table}[t]
\centering
\scriptsize
\setlength{\tabcolsep}{4pt}
\caption{External-archive statistics of MO-Fuzzy (mean\,$\pm$\,SD over folds): normalised hypervolume on inner out-of-fold objectives and on the outer test fold (refitted archive members), archive size, and modal selected $\lambda$ of the deployed rule base.}
\label{tab:pareto}
\begin{adjustbox}{max width=\linewidth}
\begin{tabular}{lcccc}
\toprule
Dataset & HV (inner) & HV (test) & Archive size & Modal $\lambda$ \\
\midrule
Pima & 0.270\,$\pm$\,0.017 & 0.264\,$\pm$\,0.073 & 10.4 & 0.1 \\
Cleveland & 0.491\,$\pm$\,0.036 & 0.429\,$\pm$\,0.104 & 34.1 & 0.1 \\
WDBC & 0.900\,$\pm$\,0.011 & 0.873\,$\pm$\,0.041 & 24.8 & 0.01 \\
ILPD & 0.171\,$\pm$\,0.010 & 0.142\,$\pm$\,0.044 & 26.6 & 0.1 \\
Thyroid & 0.880\,$\pm$\,0.015 & 0.883\,$\pm$\,0.022 & 84.7 & 0.01 \\
CKD & 0.963\,$\pm$\,0.004 & 0.948\,$\pm$\,0.025 & 38.7 & 0.01 \\
\bottomrule
\end{tabular}
\end{adjustbox}
\end{table}

Fig.~\ref{fig:pareto} shows, for every dataset, test-fold AUROC and Brier score of all
archive members against their size. Discrimination saturates after very few literals
on all datasets: further literals mainly lower the Brier score, and only slowly. The
$\varepsilon$-compact rule therefore selects models near the knee of the front. The
test-fold hypervolume of the archive was close to its inner-validation hypervolume
(Table~\ref{tab:pareto}), so the fronts found inside the training folds carried over to
unseen patients. The largest drop was on Cleveland, the smallest dataset. The
preregistered best-Brier ablation, taken from the same runs, used
\resMeanLiteralsAllMOFuzzyBestBrier{} literals on average instead of
\resMeanLiteralsAllMOFuzzyNSGAII{}. It lowered the mean Brier score from
\resMeanBrierAllMOFuzzyNSGAII{} to \resMeanBrierAllMOFuzzyBestBrier{} at an essentially
unchanged mean AUROC (\resMeanAurocAllMOFuzzyBestBrier{} against
\resMeanAurocAllMOFuzzyNSGAII{}). The front thus makes the trade explicit: the extra
literals buy a small calibration gain, and the user can choose where to stop.

\subsection{Model size and computational cost}
\label{sec:complexity}
\begin{table}[t]
\centering
\scriptsize
\setlength{\tabcolsep}{3pt}
\caption{Interpretability and cost, averaged over all datasets and folds. Rules = distinct active rules (fuzzy) or leaves (tree); literals = antecedent conditions summed over rules (tree: sum of leaf depths). Fit time includes inner tuning or evolution, single-threaded.}
\label{tab:complexity}
\begin{adjustbox}{max width=\linewidth}
\begin{tabular}{lrrrrr}
\toprule
Model & Rules & Ante./rule & Literals & Fit (s) & Peak RSS (MB) \\
\midrule
LogReg & -- & -- & -- & 1.1 & 265 \\
Decision tree & 4.8 & 2.21 & 11.4 & 1.3 & 265 \\
XGBoost & -- & -- & -- & 5.7 & 277 \\
LightGBM & -- & -- & -- & 5.1 & 275 \\
CatBoost & -- & -- & -- & 55.2 & 275 \\
MLP & -- & -- & -- & 5.5 & 273 \\
XGBoost + Platt & -- & -- & -- & 5.9 & 282 \\
XGBoost + isotonic & -- & -- & -- & 5.8 & 283 \\
SO-GA-Fuzzy (classic) & 8.1 & 2.07 & 16.6 & 11.0 & 282 \\
SO-GA-Fuzzy (hybrid) & 10.5 & 2.07 & 21.9 & 70.6 & 205 \\
MO-Fuzzy (NSGA-II) & 3.3 & 1.15 & 3.9 & 48.3 & 199 \\
MO-Fuzzy (best-Brier) & 6.3 & 1.54 & 9.8 & -- & 199 \\
\bottomrule
\end{tabular}
\end{adjustbox}
\end{table}

The deployed MO-Fuzzy models had on average \resMeanRulesAllMOFuzzyNSGAII{} rules and
\resMeanLiteralsAllMOFuzzyNSGAII{} literals (Table~\ref{tab:complexity}). The
single-objective hybrid had \resMeanRulesAllSOGAFuzzyHybrid{} rules and
\resMeanLiteralsAllSOGAFuzzyHybrid{} literals, the classic genetic-fuzzy model
\resMeanRulesAllSOGAFuzzyClassic{} rules and \resMeanLiteralsAllSOGAFuzzyClassic{}
literals; for the tuned decision tree, the table gives leaves and summed leaf depths. A
MO-Fuzzy fit, including 6{,}400 genome evaluations each with inner 3-fold
cross-validation over four penalties, took \resFitMOFuzzyNSGAII{}\,s on one CPU thread
on average, comparable to a tuned CatBoost model (\resFitCatBoost{}\,s). No fit needed
more than \resMaxRSS{}\,MB of memory.
